\documentclass[12pt]{article}

% Configuración básica
\usepackage[spanish]{babel}
\usepackage[utf8]{inputenc}
\usepackage[T1]{fontenc}
\usepackage{geometry}
\usepackage{setspace}
\usepackage{hyperref}
\usepackage{enumitem}
\usepackage{graphicx}
\usepackage[section]{placeins} 
\usepackage{float}

\geometry{letterpaper, margin=2.5cm}
\onehalfspacing

% Datos del informe (puedes modificarlos)
\title{Avance 7 \\ Participación en actividades ágiles y reuniones de seguimiento}
\author{David Armando Ramírez Navarrete}
\date{08 de febrero de 2026}

\begin{document}

\maketitle

\section*{Resumen}
El presente informe técnico documenta la participación en actividades ágiles y reuniones de
seguimiento desarrolladas durante las prácticas profesionales, como parte del trabajo
colaborativo orientado a la organización, planeación y control de actividades. El informe
describe la asistencia a reuniones periódicas de seguimiento, planeación de iteraciones y
retrospectivas, así como el desarrollo de actividades formativas relacionadas con
metodologías ágiles.

Adicionalmente, se documenta la elaboración de un sitio web estático con fines académicos
y de autoaprendizaje, utilizado como guía básica para la comprensión de los principios,
roles y eventos del marco de trabajo ágil. Finalmente, se presenta el proceso de evaluación
correspondiente a una certificación en metodologías ágiles, en el cual se obtuvo un
resultado aprobatorio.

El trabajo refleja el desarrollo de competencias en trabajo colaborativo, comunicación
efectiva, organización de actividades y aplicación de prácticas ágiles, fortaleciendo el
perfil profesional en el área de Ingeniería en Sistemas.


\section*{Introducción}
Las prácticas profesionales representan una oportunidad para aplicar conocimientos
académicos en un entorno real de trabajo, fortaleciendo competencias técnicas y
profesionales a través de la participación activa en dinámicas organizacionales. En este
contexto, la adopción de prácticas ágiles permite mejorar la organización del trabajo, la
comunicación dentro de los equipos y el seguimiento de las actividades asignadas.

Durante el desarrollo de las prácticas profesionales se participó en actividades ágiles
orientadas a la planeación, seguimiento y mejora continua del trabajo colaborativo. Dichas
actividades incluyeron la asistencia a reuniones periódicas de seguimiento, sesiones de
planeación de iteraciones y espacios de retrospectiva, así como el desarrollo de actividades
formativas complementarias relacionadas con metodologías ágiles.

El presente informe documenta dichas actividades, describiendo la metodología de trabajo
colaborativo empleada, el seguimiento de avances y los resultados obtenidos, con el
propósito de evidenciar la aplicación de prácticas ágiles y su impacto en la organización
del trabajo y en el desempeño profesional durante la estadía.


\section*{Descripción de las actividades ágiles realizadas}

Las actividades desarrolladas durante las prácticas profesionales se enmarcaron en un
entorno de trabajo colaborativo basado en prácticas ágiles, orientadas a la organización,
planeación y seguimiento de actividades. Dicho entorno involucró la participación en
dinámicas de coordinación periódica, cuyo objetivo fue asegurar la alineación del equipo,
la revisión continua de avances y la identificación oportuna de posibles impedimentos.

Como parte de estas actividades, se asistió de manera constante a reuniones de seguimiento
diario, sesiones de planeación de iteraciones y espacios de retrospectiva. Las reuniones de
seguimiento diario permitieron comunicar el estado de las tareas asignadas, identificar
bloqueos y coordinar acciones inmediatas. Por su parte, las sesiones de planeación se
enfocaron en la definición y priorización de actividades a desarrollar durante cada
iteración, considerando tiempos, dependencias y objetivos establecidos.

Asimismo, las sesiones de retrospectiva facilitaron la evaluación del desempeño del equipo
y del proceso de trabajo, promoviendo la identificación de áreas de mejora y la adopción de
acciones correctivas orientadas a la mejora continua. Estas dinámicas fortalecieron la
comunicación, la colaboración y la capacidad de adaptación ante cambios en el entorno de
trabajo.

Adicionalmente, como actividad complementaria de carácter formativo, se desarrolló un
sitio web estático con fines académicos, el cual funcionó como una guía básica de consulta
sobre los principios, roles y eventos del marco de trabajo ágil. Este recurso permitió
reforzar la comprensión conceptual de la metodología y apoyar la participación informada
en las actividades ágiles del equipo.

Finalmente, se llevó a cabo el proceso de evaluación correspondiente a una certificación
en metodologías ágiles, obteniendo un resultado aprobatorio. Este logro valida los
conocimientos adquiridos y evidencia la correcta comprensión y aplicación de las prácticas
ágiles en un contexto de trabajo colaborativo.


\section*{Metodología de trabajo colaborativo}

La metodología de trabajo empleada durante las prácticas profesionales se basó en la
aplicación de prácticas ágiles orientadas a la colaboración, la comunicación constante y la
organización estructurada de las actividades. El trabajo se desarrolló bajo un enfoque
iterativo, permitiendo planear, ejecutar y revisar las tareas de manera continua, conforme a
los objetivos definidos por el equipo.

La planeación de actividades se realizó de forma colaborativa, estableciendo prioridades y
alcances claros para cada iteración. Este proceso facilitó la distribución equilibrada de
tareas, la gestión del tiempo y la atención oportuna a dependencias entre actividades,
favoreciendo un flujo de trabajo ordenado y controlado.

El seguimiento de las actividades se apoyó en reuniones periódicas de coordinación, en las
cuales se revisó el avance de las tareas, se identificaron posibles impedimentos y se
definieron acciones correctivas cuando fue necesario. Estas dinámicas fortalecieron la
comunicación entre los integrantes del equipo y promovieron la responsabilidad compartida
sobre los resultados.

Finalmente, la metodología incorporó espacios de evaluación y retroalimentación orientados
a la mejora continua, en los cuales se analizaron los resultados obtenidos y se identificaron
oportunidades de mejora tanto a nivel individual como colectivo. Este enfoque contribuyó a
fortalecer el trabajo colaborativo y a mejorar la eficiencia en el desarrollo de las
actividades durante la estadía profesional.


\section*{Seguimiento de actividades y avances}

El seguimiento de las actividades se realizó de manera continua a lo largo del desarrollo de
las prácticas profesionales, con el objetivo de asegurar el cumplimiento de los compromisos
establecidos y mantener una visión clara del estado de avance del trabajo. Este proceso se
basó en la revisión periódica de las tareas asignadas, considerando su nivel de progreso,
prioridad y relación con los objetivos definidos para cada iteración.

Durante las reuniones de seguimiento se comunicó el avance de las actividades, se
identificaron posibles impedimentos y se acordaron acciones para su atención. Este
mecanismo permitió detectar oportunamente desviaciones respecto a lo planeado y realizar
ajustes en la organización del trabajo, favoreciendo la continuidad del flujo de actividades.

Asimismo, el seguimiento incluyó la verificación de los resultados obtenidos al cierre de
cada iteración, evaluando el grado de cumplimiento de las tareas y la calidad de los
entregables generados. La información obtenida durante este proceso sirvió como base para
la planeación de actividades posteriores y para la priorización de tareas con mayor impacto.

En conjunto, el seguimiento sistemático de actividades y avances contribuyó a mejorar la
transparencia del trabajo, fortalecer la coordinación del equipo y asegurar un control más
efectivo del desarrollo de las actividades durante la estadía profesional.


\section*{Resultados obtenidos}

Como resultado de la participación en actividades ágiles y del seguimiento continuo del
trabajo colaborativo, se lograron avances significativos tanto en el ámbito organizacional
como en el formativo. La asistencia constante a reuniones de seguimiento, planeación de
iteraciones y retrospectivas permitió fortalecer la comprensión práctica de los principios
ágiles, así como mejorar la comunicación, coordinación y alineación de objetivos dentro del
equipo de trabajo.

Uno de los principales resultados obtenidos fue el desarrollo de un sitio web estático con
fines académicos y formativos, diseñado como una guía básica para el estudio de
metodologías ágiles. Este recurso integra de manera estructurada conceptos teóricos,
eventos, roles, artefactos y valores fundamentales, además de incluir ejemplos prácticos y
un módulo de autoevaluación. El desarrollo del sitio permitió aplicar conocimientos de
estructura web, organización de contenidos y documentación técnica, al tiempo que reforzó
la comprensión conceptual del marco de trabajo ágil.

Asimismo, el uso del sitio web como herramienta de apoyo facilitó un proceso de estudio
autónomo y sistemático, permitiendo evaluar el nivel de comprensión de los contenidos y
detectar áreas de oportunidad previo al proceso de certificación. Como resultado de este
proceso formativo, se llevó a cabo la evaluación correspondiente a la certificación en
metodologías ágiles, obteniendo un resultado aprobatorio.

En conjunto, los resultados obtenidos evidencian el desarrollo de competencias en trabajo
colaborativo, organización de actividades, autoaprendizaje y aplicación de prácticas ágiles
en un entorno profesional, contribuyendo al fortalecimiento del perfil académico y
profesional en el área de Ingeniería en Sistemas.


\section*{Conclusiones}

La participación en actividades ágiles y reuniones de seguimiento durante el desarrollo de
las prácticas profesionales permitió consolidar aprendizajes relevantes relacionados con la
organización del trabajo, la comunicación efectiva y el trabajo colaborativo. La aplicación
de prácticas ágiles favoreció un enfoque estructurado y adaptable para la planeación,
ejecución y control de las actividades asignadas, contribuyendo a un mejor seguimiento de
los avances y al cumplimiento de los objetivos establecidos.

El desarrollo de un sitio web estático con fines formativos representó una experiencia
significativa de integración entre el aprendizaje teórico y la aplicación práctica. Este
recurso permitió reforzar la comprensión de los principios, roles, eventos y valores del
marco de trabajo ágil, además de fomentar el autoaprendizaje y la organización sistemática
del estudio en un contexto profesional.

Asimismo, la obtención de un resultado aprobatorio en la certificación en metodologías
ágiles valida los conocimientos adquiridos y evidencia la capacidad para aplicar prácticas
ágiles en entornos colaborativos de trabajo. Este logro fortalece el perfil académico y
profesional, al demostrar iniciativa, disciplina y compromiso con la mejora continua.

En conjunto, las actividades desarrolladas durante esta etapa de la estadía profesional
contribuyeron al fortalecimiento de competencias clave para la Ingeniería en Sistemas, tales
como la gestión de actividades, la adaptación al cambio y la colaboración efectiva, las
cuales resultan fundamentales para el desempeño en entornos profesionales actuales.


\section*{Referencias}

\begin{itemize}
    \item Google Cloud. (n.d.). \textit{Cómo crear campos calculados en Looker Studio}.  
    Recuperado de: \url{https://docs.cloud.google.com/looker/docs/studio/tutorial-create-calculated-fields-in-looker-studio?hl=es-419}

    \item Google Cloud. (n.d.). \textit{Variables de resultado de consulta en Looker Studio}.  
    Recuperado de: \url{https://docs.cloud.google.com/looker/docs/studio/query-result-variables?hl=es-419}

    \item Google Cloud. (n.d.). \textit{Agregar, editar y solucionar problemas con campos calculados en Looker Studio}.  
    Recuperado de: \url{https://docs.cloud.google.com/looker/docs/studio/add-edit-and-troubleshoot-calculated-fields?hl=es-419}
\end{itemize}

\appendix
\section{Evidencias gráficas de la Actividad 7}



\end{document}

